\documentclass[11pt]{article}
\usepackage[a4paper,margin=21mm]{geometry}
\usepackage[no-math]{fontspec}
\setmainfont{Latin Modern Roman}
\newfontfamily\CJKfont{Noto Serif CJK SC}
\XeTeXlinebreaklocale "zh"
\XeTeXlinebreakskip=0pt plus 1pt
\usepackage{amsmath,amssymb,amsthm,amscd,graphicx,color}
\usepackage{xurl}
\usepackage[unicode,hidelinks]{hyperref}
\allowdisplaybreaks
\setlength\emergencystretch{3em}
\numberwithin{equation}{section}
\numberwithin{equation}{section}

\newtheorem{Theorem}{Theorem}[section]
\newtheorem*{Theorem*}{Theorem}
\newtheorem{Corollary}[Theorem]{Corollary}
\newtheorem{Lemma}[Theorem]{Lemma}
\newtheorem{Proposition}[Theorem]{Proposition}

\theoremstyle{definition}
\newtheorem{Definition}[Theorem]{Definition}
\newtheorem{Note}[Theorem]{Note}
\newtheorem{Example}[Theorem]{Example}
\newtheorem{Remark}[Theorem]{Remark}

\usepackage{amscd,mathtools,bm}
\usepackage[all,cmtip]{xy}
\setcounter{MaxMatrixCols}{12}
\usepackage{tikz}
\usepackage{tikz-cd}
\usepackage{enumitem}
\usepackage{multirow}

\newcommand{\Section}[1]{\hyperref[sec:#1]{Section~\ref*{sec:#1}}}
\newcommand{\Table}[1]{\hyperref[tab:#1]{Table~\ref*{tab:#1}}}
\newcommand{\eqn}[1]{\hyperref[eqn:#1]{(\ref*{eqn:#1})}}
\newcommand{\Figure}[1]{\hyperref[fig:#1]{Figure~\ref*{fig:#1}}}

\newcommand{\dbb}[1]{[\![#1]\!]}
\newcommand{\br}[1]{\left(#1\right)}
\newcommand{\f}[1]{\mathfrak{#1}}

\makeatletter
\def\namedlabel#1#2{\begingroup
 #2%
 \def\@currentlabel{#2}%
 \phantomsection\label{#1}\endgroup
}
\makeatother

\DeclareMathOperator{\GL}{GL}
\DeclareMathOperator{\SL}{SL}
\DeclareMathOperator{\Sp}{Sp}
\DeclareMathOperator{\SO}{SO}
\DeclareMathOperator{\Gr}{Gr}
\DeclareMathOperator{\Fl}{Fl}
\DeclareMathOperator{\LG}{LG}
\DeclareMathOperator{\IG}{IG}
\DeclareMathOperator{\SG}{SG}
\DeclareMathOperator{\SF}{SF}
\DeclareMathOperator{\OG}{OG}
\DeclareMathOperator{\OF}{OF}
\DeclareMathOperator{\Mat}{Mat}
\DeclareMathOperator{\Hom}{Hom}
\DeclareMathOperator{\Aut}{Aut}
\DeclareMathOperator{\Ker}{Ker}
\DeclareMathOperator{\Span}{Span}
\DeclareMathOperator{\Spec}{Spec}
\DeclareMathOperator{\QH}{QH}
\DeclareMathOperator{\QK}{QK}
\DeclareMathOperator{\K}{K}
\DeclareMathOperator{\A}{A}
\DeclareMathOperator{\HH}{H}
\DeclareMathOperator{\pt}{pt}
\DeclareMathOperator{\chr}{char}
\DeclareMathOperator{\gr}{gr}
\DeclareMathOperator{\codim}{codim}
\DeclareMathOperator{\Lie}{Lie}
\DeclareMathOperator{\Pic}{Pic}
\DeclareMathOperator{\Rep}{Rep}
\DeclareMathOperator{\diag}{diag}
\DeclareMathOperator{\dist}{dist}
\def\poly{{\mathrm{poly}}}
\def\loc{{\mathrm{loc}}}
\DeclareMathOperator{\Ch}{ch}
\DeclareMathOperator{\End}{end}
\DeclareMathOperator{\supp}{supp}

\DeclareMathOperator{\ch}{\ch}
\newcommand{\Oh}{\mathcal{O}}
\DeclareMathOperator{\QKTpoly}{QK^{poly}_{\mathit{T}}}
\DeclareMathOperator{\Tor}{\textbf{\textrm{\underline{Tor}}}}
%\usepackage{soul}
%\setul{}{0.5pt}
\usepackage{bm}
\newcommand{\lbr}{\bm{[}}
\newcommand{\rbr}{\bm{]}}

\DeclarePairedDelimiter{\angles}{\langle}{\rangle}

\newcommand{\ssm}{\smallsetminus}
\newcommand{\bfr}{{\mathbf r}}
\newcommand{\bP}{{\mathbb P}}
\renewcommand{\P}{{\mathbb P}}
\newcommand{\C}{{\mathbb C}}
\newcommand{\R}{{\mathbb R}}
\newcommand{\Q}{{\mathbb Q}}
\newcommand{\Z}{{\mathbb Z}}
\newcommand{\N}{{\mathbb N}}
\newcommand{\cE}{{\mathcal E}}
\newcommand{\cF}{{\mathcal F}}
\newcommand{\cL}{{\mathcal L}}
\newcommand{\cM}{{\mathcal M}}
\newcommand{\cO}{{\mathcal O}}
\newcommand{\cP}{{\mathcal P}}
\newcommand{\cQ}{{\mathcal Q}}
\newcommand{\cR}{{\mathcal R}}
\newcommand{\cS}{{\mathcal S}}
\newcommand{\fS}{{\mathfrak S}}
\newcommand{\fQ}{{\mathfrak Q}}
\newcommand{\cT}{{\mathcal T}}
\newcommand{\cZ}{{\mathcal Z}}
\newcommand{\rmu}{{\mathrm u}}
\newcommand{\rmw}{{\mathrm w}}
\newcommand{\ds}{\displaystyle}
% \newcommand{\gw}[2]{\langle #1 \rangle^{\mbox{}}_{#2}}
% \newcommand{\egw}[2]{\langle #1 \rangle^T_{#2}}
\newcommand{\gw}[2]{\angles{ #1 }^{\mbox{}}_{#2}}
\newcommand{\egw}[2]{\angles{ #1 }^T_{#2}}
\newcommand{\euler}[1]{\chi_{_{#1}}}
\newcommand{\id}{\text{id}}
\newcommand{\bull}{{\scriptscriptstyle \bullet}}
\newcommand{\al}{{\alpha}}
\newcommand{\be}{{\beta}}
\newcommand{\ga}{{\gamma}}
\newcommand{\vep}{{\varepsilon}}
\newcommand{\ka}{{\kappa}}
\newcommand{\la}{{\lambda}}
\newcommand{\La}{{\Lambda}}

\newcommand{\Bl}{\text{B$\ell$}}
\newcommand{\sh}{\operatorname{shape}}
\DeclareMathOperator{\ev}{ev}
\DeclareMathOperator{\Jac}{Jac}
\newcommand{\wt}{\widetilde}
\newcommand{\wh}{\widehat}
\newcommand{\wb}{\overline}
\newcommand{\ov}{\overline}
\newcommand{\pic}[2]{\includegraphics[scale=#1]{#2}}
\newcommand{\ignore}[1]{}
\newcommand{\fixme}[1]{{\bf (*** #1 ***)}}
\newcommand{\Mb}{\wb{\mathcal M}}

\newcommand{\Gq}{\Gamma\llbracket q\rrbracket}
\newcommand{\noin}{\noindent}
\newcommand{\lex}{{\operatorname{lex}}}

\newcommand{\M}{\overline{M}}

\begin{document}
\begin{center}\Large Equivariant quantum-K Toda presentation for a type-A partial flag variety\end{center}
\noindent\textbf{Stable ID:} 811\quad\textbf{Author:} Kamyar Amini; Irit Huq-Kuruvilla; Leonardo C. Mihalcea; Daniel Orr; Weihong Xu\par
\noindent\textbf{Work:} Toda-Type Presentations for the Quantum K Theory of Partial Flag Varieties\par
\noindent\textbf{Source:} \url{https://sigma-journal.com/2025/098/}\par
\noindent\textbf{Version:} Official SIGMA 21 (2025) journal PDF and native TeX; acquired 2026-09-30; exact bytes pinned by SHA256\par
\noindent\textbf{Rights:} CC BY 4.0. Authors retain copyright. \url{https://creativecommons.org/licenses/by/4.0/}. Reuse and typesetting adaptation; original language and mathematical content preserved.\par
\noindent\textbf{Difficulty (prior selection preserved):} {\CJKfont H2: The author transports quantum determinant relations through partial-flag forgetting maps and proves the required pushforward identities using Grassmann-bundle geometry and the quantum corrections. The formal complete-Nakayama step is not treated as the whole proof; the nonclassical relations and their geometric construction are kept as core.}\par
\medskip\noindent\textbf{Editorial boundary note (not author text):} Complete original Theorem3.4, partial-flag/rank/tautological-bundle/torus/Novikov-completion hypotheses and all actual author determinant rewriting, localized pushforward and induction, quantum product/pushforward identities and presentation-completeness argument. The logical roadmap is retained exactly, including the authors’ explicit statement that they do not attempt self-contained proofs. The full statements and exact prior references for Kato’s homomorphism, Maeno–Naito–Sagaki’s determinant identity and the quantum Nakayama/classical Whitney prerequisites remain visible. All new current-paper pushforward and determinant-comparison calculations are retained. AppendixA’s actual Toda-from-difference-operators derivation is included because the author defers the full-flag Toda starting presentation account there. The current Whitney-to-Toda elimination and its hypotheses are included as the roadmap’s complete comparison context, supporting only. Independent later Schubert-polynomial representatives and illustrative examples are excluded. All source formulas, including relation-polynomial display conventions, are unchanged; no editor proof, hypothesis transfer or mathematical repair is supplied.\par
\medskip\hrule\medskip\noindent\textbf{Author text begins below.}\par
\setcounter{section}{1}
% Original source lines 195--216

Let $\Fl(n)$ denote the variety of complete flags in $\C^n$, and let
$\Fl(\bfr,n)=\Fl(r_1, \dots, r_k; n)$ be the variety of partial flags. These are homogeneous under the group
$\SL_n(\C)$, and the restriction of this action to the maximal torus
$T \subset \SL_n(\C)$ has finitely many fixed points, indexed
by a~quotient of the symmetric group $S_n$.
Denote
by $\QK_T(\Fl(\bfr,n))$ the (equivariant, small) quantum K
ring associated to these varieties. This is an algebra over
$\K_T(\pt)[\![Q_1, \dots, Q_k]\!]$, and it has
a~$\K_T(\pt)[\![Q_1, \dots, Q_k]\!]$-basis given by Schubert classes
$\cO^w$ indexed by the torus fixed points. The quantum K multiplication
was defined by Givental and Lee~\cite{givental:onwdvv,lee:QK} in
terms of $3$-point, genus~$0$, $\K$-theoretic Gromov--Witten (KGW) invariants.
Denote by
\[ 0=\cS_0 \subset \cS_1 \subset \dots \subset \cS_k \subset \cS_{k+1}=\C^n \]
the sequence of tautological bundles in $\Fl(r_1, \dots, r_k; n)$; thus
$\mathrm{rank}(\cS_i)=r_i$ for $0\leq i\leq k+1$ with $r_0=0$ and $r_{k+1}=n$.

While the computational foundations of the quantum K rings
of (cominuscule) Grassmannians have been studied for some
time now (see, e.g.,~\cite{BCMP:qkpos,BCMP:qkchev,buch.m:qk,chaput.perrin:rationality,Gorbounov:2014,sinha2024quantum}),


% Original source lines 237--248
and also by Anderson, Chen and Tseng in~\cite{act:2017}, see also
\cite{koroteev} and Appendix~\ref{sec:toda} below.

Our main result is to generalize the Toda presentation from $\QK_T(\Fl(n))$
to one for the ring
$\QK_T(\Fl(\bfr,n))$ associated to partial flag varieties. To state it, let
\[
Y^{(j)}=\bigl(Y^{(j)}_1,\dots,{ Y^{(j)}_{r_{j+1}-r_j}}\bigr), \qquad 0\leq j\leq k
\]
be formal variables and $e_\ell$ be the $\ell$-th elementary symmetric polynomial. Let $T_1,\dots, T_n \in \K_T(\pt)$ be given by
the decomposition of $\C^n$ into one dimensional $T$-modules, that is, $\wedge^\ell(\C^n)=e_\ell(T_1,\dots,T_n)$. To distinguish from multiplication in $K_T(\Fl(\bfr,n))$, we denote the multiplication in $\QK_T(\Fl(\bfr,n))$ by $\star$.



% Original source lines 274--290
Our proof is by decreasing induction on $k$. The initial case, $k=n-1$, is the main result of~\cite{maeno.naito.sagaki:QKideal}, rewritten in determinantal form in Theorem~\ref{thm:toda} below. For the induction step, we use a~result of Kato~\cite{kato2019quantum} which states that there is a
$\K_T(\pt)$-{\em algebra} homomorphism
\begin{gather}
\QK_T(\Fl(n)) \to \QK_T(\Fl(r_1, \dots , r_k;n)); \nonumber\\
\qquad{}\cO^w \mapsto \pi_*(\cO^w) ,
\qquad Q_r \mapsto \begin{cases} 1, & r \notin \{ r_1, \dots, r_k \}, \\
Q_i, & r=r_i ,\end{cases}\label{E:kato-push}
\end{gather}
which extends the usual projection map $\pi_*\colon \K_T(\Fl(n)) \to \K_T(\Fl(\bfr,n))$.
Note that the classical~$\pi_*$ is {\em not} a ring map. (Kato's result is for general complex, simple groups $G$.) We use this to show that the ideal $J_Q$ is contained in the ideal of relations.
For the specialization ${Q_i \mapsto 1}$ to be well defined, one needs to work with
polynomials in $Q_1, \dots, Q_{n-1}$; see Section~\ref{sec:prelim}.
Pushing forward the original Toda relations is not possible, due to poles at $Q_i=1$.
We had to rewrite these relations, and additionally
use an extra identity due to Maeno, Naito, and Sagaki
(cf.~Proposition~\ref{prop:MNS} below), in order for the push forward to be performed.
Finally, it follows from~\cite{GMSXZZ:Nakayama} that the ideal $J_Q$ coincides with the ideal of relations.


% Original source lines 358--854
\subsection{A logical roadmap} There are several recent results in the literature which inform the Toda and Whitney
presentations we prove in this paper.
Since we do not attempt to give self-contained
proofs, we provide next a roadmap of the logical
implications we rely on, which a reader may find useful.
 %(Each of these will be recalled again during the relevant proofs.)

Our proofs of both the Toda presentation, and the Whitney presentation,
for $\QK_T(\Fl(\bfr,n))$, from Theorem~\ref{thm:main} (resp.\ \eqref{E:QKWhitney}, see Remark~\ref{rmk:QKW})
rely on the following: the Toda (resp.\ Whitney) presentation for $\QK_T(\Fl(n))$;
Kato's push forward homomorphism from~\eqref{E:kato-push}; and the key
technical result from Proposition~\ref{prop:MNS}, proved in~\cite{maeno.naito.sagaki:QKideal},
which in turn relies on Kato's work~\cite{kato:loop}.

There are two proofs of the Toda presentation for $\QK_T(\Fl(n))$, one in
\cite{maeno.naito.sagaki:QKideal}, relying on~\cite{kato:loop}, and another
which may be deduced from~\cite{act:2017}, relying on results from
\cite{givental_lee} and~\cite{Iritani:2013qka};~cf.~Appendix~\ref{sec:toda}.

There are also two proofs of the Whitney presentation for $\QK_T(\Fl(n))$.
One was recently announced by Huq-Kuruvilla
\cite{huq2024quantum} (for all rings $\QK_T(\Fl(\bfr,n))$), and
it uses the technique of abelian-nonabelian correspondence, independent
of Kato's results. Another proof of the Whitney presentation for
 $\QK_T(\Fl(n))$ is given in~\cite[Section~6]{GMSXZZ:QKW}.
It relies on the recent proof of the quantum~K divisor axiom~\cite{LNSX:QKdiv},
and ultimately on Kato's results.

\section{Preliminaries}
\subsection{Equivariant K theory of Grassman bundles}
Let $T$ be a linear algebraic group. For any projective $T$-variety $Z$,
let $\K_T(Z)$ be the equivariant $\K$-theory
ring, defined as the Grothendieck ring of $T$-equivariant algebraic vector
bundles. This ring is an algebra over $\K_T(\pt)$, the representation ring of
$T$. Let $\euler{Z}\colon \K_T(Z) \to \K_T(\pt)$ be the pushforward map along the
structure morphism.

For $E \to Z$ a $T$-equivariant
vector bundle of rank $\operatorname{rk}E$, we denote by
\[ \lambda_y(E) := 1 + y [E] + \dots + y^{\operatorname{rk}E} \bigl[\wedge^{\operatorname{rk}E} E\bigr] \in \K_T(Z)[y] \]
the Hirzerbruch $\lambda_y$ class of $E$. This class is multiplicative for short exact sequences. {In an abuse of notation, we often write $E$ for the class $[E]$ in $\K_T(Z)$.} Note that for a rank $e$ equivariant vector bundle $E$, and a character ${\rm e}^\chi \in \K_T(\pt)$,
\[ \lambda_y( {\rm e}^\chi \otimes E) = \lambda_{y {\rm e}^{\chi}}(E) = \sum_{i=0}^{e} y^i {\rm e}^{i \chi} \otimes \wedge^i E . \]
As is customary, we will often remove the $\otimes$ symbol from the notation.

Denote by
$\pi\colon \mathbb{G}(r,E) \to Z$ the Grassmann bundle over $Z$.
It is equipped with a tautological sequence
$0 \to \underline{\cS} \to \pi^* E \to \underline{\cQ} \to 0$ over
$\mathbb{G}(r,E)$. The following result follows
from~\cite[Proposition~2.2]{kapranov:Gr}, see also~\cite[Proposition~3.2 and Corollary~3.3]{gu2022quantum}.
(Kapranov proved this when $Z=\pt$; the relative version follows immediately
using that $\pi$ is a $T$-equivariant locally trivial fibration.) We only state the special cases that will be used in this paper. See the above references for the full generality.

\begin{Proposition}[Kapranov]\label{prop:relativeBWB}
 There are the following isomorphisms of $T$-equivariant vector bundles:
 \begin{enumerate}\itemsep=0pt
 \item[$(1)$] for all $i \ge 0$, $\ell>0$ the higher direct images, $R^i \pi_* \bigl(\wedge^\ell\underline{\cS}\bigr) = 0$;
 \item[$(2)$] for all $\ell\geq 0$, \[ R^i \pi_*\bigl(\wedge^\ell\underline{\cQ}\bigr) = \begin{cases} \wedge^\ell E, & i=0, \\ 0, & i>0 . \end{cases}\]
 \end{enumerate}
 \end{Proposition}

\subsection{(Equivariant) quantum K theory of flag varieties}\label{sec:prelim}
Let $\mathbf{r}= (r_1, \dots, r_k)$. We consider \[X=\Fl(\bfr,n),\] which parametrizes flags of vector spaces
$F_{1} \subset F_{2} \subset \dots \subset F_{k} \subset \C^n$
with $\dim F_{i} = r_i$ for~${1 \le i \le k}$.

Let $M_{d,n} := \Mb_{0,n}(X,d)$ be the moduli space of genus zero degree
$d$ stable maps to $X$ with~$n$ marked points. Given classes
$a_1, \dots, a_n \in \K_T(X)$,
define the $\K$-theoretic Gromov--Witten invariants by
\[ \langle a_1, \dots, a_n \rangle_{d} =
\euler{M_{d,n}}\Biggl(\prod_{i = 1}^n\ev_i^*(a_i)\Biggr)\in \K_T(\pt) . \]
Non-equivariant Gromov--Witten invariants are obtained by replacing $T$ with the
trivial group; these Gromov--Witten invariants are integers.

For $d=(d_1,\dots,d_k)\in\HH_2(X,\Z)\cong\Z^k$, we write $Q^d$ for \smash{$\prod_{i=1}^k Q_i^{d_i}$}. Here $Q_i$ corresponds to the Poincar{\'e} dual of the first Chern class $-c_1(\det \cS_i)$. Following~\cite{givental:onwdvv,lee:QK}, the $T$-equivariant (small) quantum
K theory ring is
\begin{equation*}%\label{eqn:QKTfree}
\QK_T(X) = \K_T(X) \otimes_{\K_T(\pt)} \K_T(\pt)[\![Q]\!]
\end{equation*}
as a $\K_T(\pt)[\![Q]\!]$-module.
It is equipped with a commutative, associative
product, denoted by $\star$, which is
determined by the condition
\begin{equation*}%\label{eqn:def}
 (\!(\sigma_1\star\sigma_2,\sigma_3)\!)=\sum_d
 Q^d \gw{\sigma_1,\sigma_2,\sigma_3}{d}\qquad\text{for all $\sigma_1,\sigma_2,\sigma_3\in \K_T(X)$},
\end{equation*}
where
\[
 (\!(\sigma_1,\sigma_2)\!)\coloneqq \sum_{d}Q^d \gw{\sigma_1,\sigma_2}{d}
\]
is the quantum $\K$-metric.

It was proved in~\cite{anderson2022finite, kato:loop} that for $\sigma_1,\sigma_2\in\K_T(X)$, the product $\sigma_1\star\sigma_2$ can always be expressed as a polynomial in $Q$ with coefficients in $\K_T(X)$. It follows that
\[
 \QK_T^\poly(X)\coloneqq\K_T(X) \otimes_{\K_T(\pt)}\K_T(\pt)[Q]
\]
is a subring of $\QK_T(X)$.

Let $Y=\Fl\bigl(r_1,\dots, \widehat{r}_i,\dots, r_k;n\bigr)$ and $\pi\colon X\to Y$ be the natural map. Let also $\widehat{\mathbf{r}} = \bigl(r_1, \dots,\allowbreak \widehat{r}_i,\dots, r_k\bigr)$. The following theorem is a specialization of results proved in~\cite{kato2019quantum}.

\begin{Theorem}[Kato]\label{thm:kato-morphism}
 There is a surjective ring homomorphism
 \[
 \Phi\colon\ \QK^\poly_T(X)\to\QK^\poly_T(Y)
 \]
 given by $\sigma\mapsto{\pi}_*\sigma$ for all $\sigma\in\K_T(X)$ and
 \[
 Q_j\mapsto\begin{cases}
 Q_j, & j\neq i,\\
 1, & j=i
 \end{cases} \qquad \text{for $1\leq i\leq k$}.\]
\end{Theorem}
It follows from Theorem~\ref{thm:kato-morphism} that
Kato's homomorphism extends naturally to
\begin{equation}\label{E:Kato-loc-hom}
\Phi\colon \QK_T^{\loc(\mathbf{\hat r})}(X) \to
\QK_T^{\loc(\mathbf{\hat r})}(Y) ,
\end{equation}
where $\loc(\mathbf{\hat r})$ indicates localization at the multiplicative set generated by
$1- Q_j$ for $j \neq i$.

\subsection[The Toda presentation for Fl(n)]{The Toda presentation for $\boldsymbol{\Fl(n)}$}

The variety $\Fl(n)=\Fl(1,\dots,n-1;n)$ is equipped with tautological vector bundles
\[
 0=\cS_{0}\subset\cS_{1}\subset\dots \subset \cS_{n-1}\subset \cS_{n}=\C^n ,
\]
where $\cS_j$ has rank $r_j$. It can also be viewed as $\SL_n/B$, where $B\subset \SL_n$ is a Borel subgroup. Let~${T\subseteq B}$ be a maximal torus in $\SL_n$.

The following is the main result of~\cite{maeno.naito.sagaki:QKideal} (see Remark~\ref{rmk:mns} for more details). The relation~\eqref{eqn:toda} can also be recovered from the connection between the $J$-function of the full flag variety and the relativistic Toda lattice established by Givental and Lee in~\cite{givental_lee}. This observation was made in the unpublished note~\cite{act:2017} of Anderson--Chen--Tseng, but removed from the published version of their paper. For the sake of completeness, we give a brief account in Appendix~\ref{sec:toda}.

\begin{Theorem}\label{thm:toda}
The ring $\QK_T(\Fl(n))$ is isomorphic to $R'[\![Q]\!]/J'_Q$, where $R'$ is equal to the ring \smash{$\K_T(\pt)\bigl[P_1^{\pm },\dots,P_{n}^\pm\bigr]$} and the ideal $J'_Q\subset R'\dbb{Q}=R'\dbb{Q_1,\dots,Q_{n-1}}$ is generated by the coefficients $y$ in
 \begin{gather}\label{eqn:toda}
 \lambda_y(\C^n) -
 \begin{vmatrix}
 1+y\frac{P_1}{P_0} & y\frac{P_2}{P_1}Q_1 & & & & \\[0.9mm]
 1 & 1+y\frac{P_2}{P_1} & y\frac{P_3}{P_2}Q_2 & & & \\[0.9mm]
 & 1 & 1+y\frac{P_3}{P_2} & y\frac{P_{4}}{P_{3}}Q_3 & & \\
 & & \ddots & \ddots & \ddots & \\
& & & 1 & 1+y\frac{P_{n-1}}{P_{n-2}} & y\frac{P_n}{P_{n-1}}Q_{n-1}\\[1mm]
& & & & 1 & 1+y\frac{P_n}{P_{n-1}}
\end{vmatrix}^\star,
\end{gather}
here $P_0=1$ by convention, and $\lambda_y(\C^n)\in\K_T(\pt)[y]$.

More precisely,
there exists a $\K_T(\pt)[\![Q]\!]$-algebra isomorphism $\Psi'\colon R'[\![Q]\!]/J'_Q \to \QK_T(\Fl(n))$ that sends $P_j$ to $\det \cS_j$ for all $j=1,\dots,n$.
\end{Theorem}

\begin{Remark} \label{rmk:mns}
 Theorem~\ref{thm:toda} is proved in~\cite{maeno.naito.sagaki:QKideal} using results of Kato~\cite{kato:loop} based on the semi-infinite
 flag variety.
 		The connection between our statement of Theorem~\ref{thm:toda} and that of~\cite{maeno.naito.sagaki:QKideal} is seen as follows. Define the \textit{Toda polynomials} \smash{$T^{(n)}_k$} for $k=1,\dots,n$ by
		\begin{align*}
		T^{(n)}_k = \sum_{0= i_0<\dots<i_k\leq n}\ \prod_{s=1}^{k}\frac{P_{i_s}}{P_{i_{s}-1}}\left(1-Q_{i_s-1}\right)^{1-\delta_{i_s-i_{s-1},1}}.
		\end{align*}
		These elements of $\Z\bigl[P_1^{\pm},\dots,P_{n}^{\pm}\bigr][\![Q]\!]$ (where $P_0=1$ and $Q_0=0$ by convention) are symbols of the finite-difference Toda Hamiltonians~\cite{etingof} (see also~\cite{act:2017,gloI,givental_lee,koroteev}).
		
		Letting \smash{$T^{(n)}=\sum_{k=0}^n T^{(n)}_k y^k$} where \smash{$T^{(n)}_0=1$}, we claim that $T^{(n)}(y)$ is equal to the determinant of the matrix appearing in the Toda relations~\eqref{eqn:toda}, namely,
		\begin{align*}
 T^{(n)}(y)=\begin{vmatrix}
 1+y\frac{P_1}{P_0} & y\frac{P_2}{P_1}Q_1 & & & & \\[0.9mm]
 1 & 1+y\frac{P_2}{P_1} & y\frac{P_3}{P_2}Q_2 & & & \\[0.9mm]
 & 1 & 1+y\frac{P_3}{P_2} & y\frac{P_{4}}{P_{3}}Q_3 & & \\
 & & \ddots & \ddots & \ddots & \\
& & & 1 & 1+y\frac{P_{n-1}}{P_{n-2}} & y\frac{P_n}{P_{n-1}}Q_{n-1}\\[1mm]
& & & & 1 & 1+y\frac{P_n}{P_{n-1}}
\end{vmatrix}^\star.
		\end{align*}
		This is verified by showing that $T^{(n)}$ satisfies the recursion
		\begin{equation*}%\label{E:W-Trec}
T^{(n)}= T^{(n-1)} \biggl(1+ y \frac{P_n}{P_{n-1}}\biggr) - y Q_{n-1} \frac{P_n}{P_{n-1}} T^{(n-2)} ,
		\end{equation*}
		and then applying Lemma~\ref{L:r-det} below (with $n$ playing the role of $j$ there).
		\end{Remark}

 \begin{Remark}
 Upon the specialization $Q_1=\dotsm=Q_{n-1}=0$, the Toda presentation $R'[\![Q]\!]/J'_Q\cong\QK_T(\Fl(n))$ becomes the Borel presentation
 \[
 \K_T(\pt)\bigl[P_1^\pm,\dots,P_n^{\pm}\bigr]/J\cong \K_T(\Fl(n)),
 \]
 where $J$ is the ideal generated by the coefficients of $y$ in
 \[
 \lambda_y(\C^n)-\sideset{}{^\star}\prod_{j=0}^{n-1} (1+yP_{j+1}/P_j)
 \]
 and~$P_j$ corresponds to $\det \cS_j$ for all $j=1,\dots,n$.
 \end{Remark}
 \begin{Lemma}
 \label{L:r-det}
 Suppose $U_j$ for $0\le j\le k+1$ and $A_j$, $B_j$ for $0\le j\le k$ are elements of a~commutative ring with $1$ such that the $U_j$ satisfy the recursion
 \begin{equation*}
 %\label{E:recursion}
 U_{j+1}=A_jU_{j}-B_jU_{j-1},\qquad 0\le j\le k
 \end{equation*}
 with initial conditions $U_0=1$, $U_{-1}=0$. Then, for all $0\le j\le k+1$, one has
 \begin{equation}
 \label{E:r-det}
 U_j=
 \begin{vmatrix}
 A_0 & B_1 & & & \\
 1 & A_1 & B_2 & & \\
 & \ddots & \ddots & \ddots & \\
 & & 1 & A_{j-2} & B_{j-1}\\
 & & & 1 & A_{j-1}
 \end{vmatrix}.
 \end{equation}
 \end{Lemma}

 \begin{proof}
 One simply expands along the last row or column to see that the determinant in~\eqref{E:r-det} satisfies the recursion. Observe that the initial values $U_0=1$ and $U_1=A_0$ agree. This completes the proof.
 \end{proof}

 Before finishing this section, we record the following, which follows from~\cite[Proposition~5.2]{maeno.naito.sagaki:QKideal}.

 \begin{Proposition}[Maeno--Naito--Sagaki]\label{prop:MNS}
 In $\QK_T(\Fl(n))$, the following relations hold:
 \begin{gather*}
 \det \cS_i \star \det \cS_j/\cS_i = (1-Q_i) \det \cS_j, \qquad 1 \le i < j \le n.
 \end{gather*}
 \end{Proposition}

\section[Toda-type presentations for the equivariant quantum K theory of partial flag varieties]{Toda-type presentations for the equivariant quantum\\ K theory of partial flag varieties}%\label{sec:partial}

To begin, we observe that the Toda presentation in Theorem~\ref{thm:toda} can be rewritten as follows.
\begin{Corollary}\label{cor:fln}
 The ring $\QK_T(\Fl(n))$ is isomorphic to $R\dbb{Q}/J_Q$, where
 \[
 R=\K_T(\pt)\bigl[Y^{(0)}, \dots, Y^{(n-1)}\bigr],
 \]
 and $J_Q\subset R\dbb{Q}=R\dbb{Q_1,\dots,Q_{n-1}}$ is generated by the coefficients of $y$ in
 \begin{gather*}%\label{eqn:toda-gen}
 \lambda_y(\C^n) - \\
 \begin{vmatrix}
 1+yY^{(0)}\frac{1}{1-Q_0} & yY^{(1)}\frac{Q_1}{1-Q_1} & & & & \\[1mm]
 1 & 1+yY^{(1)}\frac{1}{1-Q_1} & yY^{(2)}\frac{Q_2}{1-Q_2} & & & \\ % & 1 & 1+yY^{(2)}\frac{1}{1-Q_2} & yY^{(3)}\frac{Q_3}{1-Q_3}Q_3 & & \\
& \ddots & \ddots & \ddots \\ & & 1 & 1+yY^{(n-2)}\frac{1}{1-Q_{n-2}} & yY^{(n-1)}\frac{Q_{n-1}}{1-Q_{n-1}}\\[1mm] & & & 1 & 1+yY^{(n-1)}\frac{1}{1-Q_{n-1}}
\!\!\!\end{vmatrix}^\star
\end{gather*}
with the convention that $Q_0=0$.

More precisely,
% for $\mathbf{r}=(1,2,\dots,n-1)$,
there exists a $\K_T(\pt)[\![Q]\!]$-algebra isomorphism $\Psi\colon R[\![Q]\!]/J_Q \to \QK_T(\Fl(n))$ that sends $Y^{(j)}$ to $\cS_{j+1}/\cS_j$ for $j=1,\dots,n-1$.
\end{Corollary}
\begin{proof}
 Identifying $P_{j+1}/P_j$ with $Y^{(j)}/(1-Q_j)$ gives an isomorphism between $R[\![Q]\!]/J_Q$ and $R'[\![Q]\!]/J'_Q$. More precisely, define a $\K_T(\pt)[\![Q]\!]$ homomorphism $\Phi\colon R'[\![Q]\!]/J'_Q\to R[\![Q]\!]/J_Q$ by
\[
 \Phi(P_j)=\prod_{i=0}^{j-1}\frac{Y^{(i)}}{1-Q_i}, \qquad 1\leq j\leq n-1.
\]
Note that in $R[\![Q]\!]/J_Q$, we have \smash{$\det\C^{n}=\prod_{i=0}^{n-1}Y^{(i)}/(1-Q_i)$}, which implies all $Y^{(j)}$ are invertible. Since the relations match, the homomorphism $\Psi$ is well-defined and injective. Since \smash{$(1-Q_j)P_{j+1}/P_j$} is sent to $Y^{(j)}$ for $0\leq j\leq n-1$, it is also surjective. Finally, the~geometric interpretation follows from Proposition~\ref{prop:MNS}.
\end{proof}

Next, we generalize Corollary~\ref{cor:fln} to all partial flag varieties utilizing Theorem~\ref{thm:kato-morphism}, Proposition~\ref{prop:MNS}, and the Nakayama-type result from~\cite{GMSXZZ:QKW, gu2022quantum}.

\begin{Theorem}\label{thm:whit-short}
 In $\QK_T(\Fl(\bfr,n))[y]$, the following relation hold:
 \begin{align}\label{eqn:whit-short}
 \lambda_y(\C^n)-
 &\begin{vmatrix}
 A_0 & B_1 & & & \\
 1 & A_1 & B_2 & & \\ & \ddots & \ddots & \ddots & \\ & & 1 & A_{k-1} & B_{k}\\ & & & 1 & A_{k}
 \end{vmatrix}^\star,
 \end{align}
 where \[
 B_j=y^{r_{j+1}-r_j}\frac{Q_j}{1-Q_j}\det(\cS_{j+1}/\cS_j),\qquad
 A_j=\lambda_y(\cS_{j+1}/\cS_j)+B_j.
\]
\end{Theorem}
\begin{proof}
Let $X=\Fl(r_1,\dots,r_k;n)$, $Y=\Fl\bigl(r_1,\dots, \widehat{r}_i,\dots, r_k;n\bigr)$, and $\pi\colon X\to Y$ be the natural map. Let
\[
 0=\cS_0\subset\cS_1\subset\dots\subset\cS_k\subset\cS_{k+1}=\C^n
\]
be the sequence of tautological bundles on $X$. Note that all but $\cS_i$ are pulled back from $Y$. With a slight abuse of notation, we denote the sequence of tautological bundles on $Y$ by
\[
 0=\cS_0\subset\cS_1\subset\dots\cS_{i-1}\subset\cS_{i+1}\subset\dots\subset\cS_k\subset\cS_{k+1}=\C^n.
\]
Note that the elements $B_1,\dots,B_{i-2},B_{i+1},\dots,B_k$ as well as $A_1.\dots, A_{i-2},A_{i+1},\dots,A_k$ in the ring $\QK_T(X)[y]$ stay the same under pushforward along $\pi$. By a slight abuse of notation, we~also think of them as elements of $\QK_T(Y)[y]$.

By induction, we assume that relation~\eqref{eqn:whit-short} holds for $X$, i.e.,
\begin{gather}\label{eqn:whitX}
 \lambda_y(\C^n)-\begin{vmatrix}
 A_0 & B_1 & & & & & & & & & \\
 1 & A_1 &B_2 & & & & & & & & \\
 & \ddots & \ddots & \ddots & & & & & & & \\
 & & \ddots & \ddots & B_{i-2} & & &\\ & & & 1 & A_{i-2} & B_{i-1} & & & & &\\ & & & & 1 & A_{i-1} & B_i & & & &\\
 & & & & & 1 & A_i & B_{i+1} & & &\\
 & & & & & & 1 & A_{i+1} & \ddots & &
 \\
 & & & & & & & \ddots & \ddots & \ddots & \\
 & & & & & & & & 1 & A_{k-1} & B_{k}\\
 & & & & & & & & & 1 & A_{k}
 \end{vmatrix}^\star
\end{gather}
holds in \smash{$\QK_T^{\loc(\mathbf{r})}(X)[y]$}
for $1 \le j \le k$, and we will show that the (localized) Kato's pushforward~\eqref{E:Kato-loc-hom} of this relation gives relation~\eqref{eqn:whit-short} on $Y$.

Relation~\eqref{eqn:whit-short} on $Y$ reads
\begin{gather}\label{eqn:whitY}
 \lambda_y(\C^n)-\begin{vmatrix}
 A_0 & B_1 & & & \\
 1 & A_1 & B_2 & & \\ & \ddots & \ddots & \ddots & \\
 & & \ddots & \ddots & B_{i-2}\\ & & & 1 & A_{i-2} & B'_{i-1}\\ & & & & 1 & A'_{i-1} & B_{i+1}\\ & & & & & 1 & A_{i+1} & \ddots
 \\ & & & & & & \ddots & \ddots & \ddots \\ & & & & & & & 1 & A_{k-1} & B_{k}\\ & & & & & & & & 1 & A_{k}
 \end{vmatrix}^\star,
\end{gather}
where
\begin{gather*}
 B_{i-1}'=y^{r_{i+1}-r_{i-1}}\frac{Q_{i-1}}{1-Q_{i-1}}\det\left(\cS_{i+1}/\cS_{i-1}\right),\qquad A_{i-1}'=\lambda_y(\cS_{i+1}/\cS_{i-1})+B'_{i-1},
\end{gather*}
regarded as elements in \smash{$\QK_T^{\loc(\widehat{\mathbf{r}})}(Y)[y]$}.

By the projection formula, to prove~\eqref{eqn:whitY}, it suffices to prove the pushforward along $\pi$ of~\eqref{eqn:whitX} agrees with~\eqref{eqn:whitY}. We compare the two determinants by expanding along columns. Expanding along the column containing $B'_{i-1}$, we have that the determinant in~\eqref{eqn:whitY} is of the form
\begin{gather*}%\label{eqn:exprX}
 -B'_{i-1}\star C'+A'_{i-1}\star D'-E';
\end{gather*}
expanding along the two columns containing $B_{i-1}$ or $B_i$, we have that the determinant in~\eqref{eqn:whitX} is of the form
\begin{gather*}
 \begin{vmatrix}B_{i-1} & 0\\
 A_{i-1} & B_{i}\end{vmatrix}^\star\star0-\begin{vmatrix}
 B_{i-1} & 0\\
 1 & A_{i}
 \end{vmatrix}^\star\star C +\begin{vmatrix}
 B_{i-1} & 0\\
 0 & 1
 \end{vmatrix}^\star\star F\\
 \qquad{} +\begin{vmatrix}
 A_{i-1} & B_i\\1 & A_i
 \end{vmatrix}^\star\star D -\begin{vmatrix}
 A_{i-1} & B_i\\
 0 & 1
 \end{vmatrix}^\star\star E+\begin{vmatrix}
 1 & A_i\\
 0 & 1
 \end{vmatrix}^\star\star 0.
 \end{gather*}
Note that $C$, $D$, $E$, $F$ stay the same under the pushforward, and it is straightforward to check that
\[
C'=C,\qquad D'=D,\qquad E'=E.
\]
The rest follows from Lemma~\ref{lemma:main-push} below.
\end{proof}

\begin{Lemma}\label{lemma:main-push} The following hold:
 \begin{enumerate}[label=$(\alph*)$]\itemsep=0pt
 \item $\pi_* \begin{vmatrix}
 A_{i-1} & B_i\\
 0 & 1
 \end{vmatrix}^\star= 1$;
 \item $\pi_* \begin{vmatrix}
 B_{i-1} & 0\\
 0 & 1
 \end{vmatrix}^\star = 0$;
 \item $\pi_* \begin{vmatrix} B_{i-1} & 0\\ 1 & A_{i} \end{vmatrix}^\star = B_{i-1}'$;
 \item Assume that $r_i - r_{i-1} = 1$. Then
 \[
 \pi_* \begin{vmatrix} A_{i-1} & B_i\\1 & A_i
 \end{vmatrix}^\star = A_{i-1}'.
 \]
 \end{enumerate}
 \end{Lemma}
 \begin{proof}
 Note that $X$ may be realized as the Grassmann bundle $\mathbb{G}(r_i-r_{i-1}, \cS_{i+1}/\cS_{i-1})$
 over~$Y$, with tautological sequence $0 \to \cS_i/\cS_{i-1} \to \cS_{i+1}/\cS_{i-1} \to \cS_{i+1}/\cS_{i} \to 0$.
 It follows from Proposition~\ref{prop:relativeBWB} that
 \begin{equation}\label{eqn:pushforwards-lambda}
 \pi_*(\lambda_y(\cS_{i+1}/\cS_{i}))=\sum_{j=0}^{r_{i+1}-r_i}y^j\wedge^j(\cS_{i+1}/\cS_{i-1}), \qquad\pi_* (\lambda_y(\cS_i/\cS_{i-1})) = 1.
 \end{equation}

 For (a), (b), note that \smash{$A_{i-1}, B_{i-1}\in \QK_T^{\loc(\mathbf{\hat r})}(X)$}, so we may use~\eqref{E:Kato-loc-hom}, and it follows that
 \begin{gather*}%\label{eqn:pushforwards}
 \pi_*B_{i-1}=0,\qquad \pi_*A_{i-1} = 1.
 \end{gather*}
 Note that by Proposition~\ref{prop:MNS} and Theorem~\ref{thm:kato-morphism}, we have
 \begin{gather}
 \det\cS_{j}\star\det(\cS_{j+1}/\cS_{j})=(1-Q_{j})\det\cS_{j+1}\qquad \text{for $0\leq j\leq k$, in $\QK_T(X)$},\label{eqn:relX}\\
 \det\cS_{i-1}\star\det(\cS_{i+1}/\cS_{i-1})=(1-Q_{i-1})\det\cS_{i+1}\qquad \text{in $\QK_T(Y)$}.\label{eqn:relY}
\end{gather}
 To prove (c), we obtain from definition
 \begin{align}
 \begin{vmatrix}
 B_{i-1} & 0\\
 1 & A_{i}
 \end{vmatrix}^\star
 &{}= B_{i-1}A_i=B_{i-1} \star (\lambda_y(\cS_{i+1}/\cS_i)+B_i)\nonumber\\
 &{}= B_{i-1}\star \lambda_y(\cS_{i+1}/\cS_i)+ B_{i-1} \star B_i .\qquad\label{eqn:3-expand}
 \end{align}
 The element $B_{i}$ cannot be pushed forward, as it contains $1-Q_i$ in the denominator. However,
 we use~\eqref{eqn:relX} to calculate
 \begin{align*}
 B_{i-1} \star B_i &{}= y^{r_{i+1}-r_{i-1}} \frac{Q_{i-1} Q_i}{(1-Q_{i-1})(1-Q_i)} \det (\cS_{i+1}/\cS_i) \star \det (\cS_{i}/\cS_{i-1}) \\
 &{}= y^{r_{i+1}-r_{i-1}}Q_{i-1}Q_i\frac{\det\cS_{i+1}}{\det\cS_{i-1}} ,
 \end{align*}
 where the inverse is calculated in the quantum K ring of $X$. By~\eqref{eqn:relX} again,
 \[ \frac{\det\cS_{i+1}}{\det\cS_{i-1}} = \frac{\det \cS_{i+1} \star\det \C^n/\cS_{i-1}}{(1-Q_{i-1})\det \C^n}\qquad \text{in $\QK_T^{\loc(\mathbf{\hat r})}(X)$}, \]
 and its pushforward is
\begin{gather}\label{eqn:3}
 \frac{\det\cS_{i+1}}{\det\cS_{i-1}}\in \QK_T^{\loc(\mathbf{\hat r})}(Y).
\end{gather}
Note that by~\eqref{eqn:relY},
expression \eqref{eqn:3} is equal to
\begin{gather*}
 \frac{\det\left(\cS_{i+1}/\cS_{i-1}\right)}{1-Q_{i-1}}\qquad \text{in} \ \QK_T^{\loc(\mathbf{\hat r})}(Y).
\end{gather*}
Using \eqref{eqn:pushforwards-lambda} and \eqref{eqn:3-expand}, the projection formula, and Theorem~\ref{thm:kato-morphism}, it follows that
\begin{gather*}
 \pi_*\begin{vmatrix}
 B_{i-1} & 0\\
 1 & A_{i}
 \end{vmatrix}^\star=\pi_*\left(B_{i-1}\star B_i\right)=y^{r_{i+1}-r_{i-1}}\frac{Q_{i-1}}{1-Q_{i-1}}\det\left(\cS_{i+1}/\cS_{i-1}\right)=B'_{i-1}.
\end{gather*}
 For (d), we calculate
 \[ \begin{vmatrix}
 A_{i-1} & B_i\\1 & A_i
 \end{vmatrix}^\star
 = \begin{vmatrix}
 \lambda_y(\cS_{i}/\cS_{i-1}) + B_{i-1} & B_i\\1 & A_i
 \end{vmatrix}^\star
 = A_i \star \lambda_y(\cS_{i}/\cS_{i-1}) + A_i \star B_{i-1} - B_i. \]
 From (c), $\pi_* (A_i \star B_{i-1}) = B_{i-1}'$, therefore it suffices to show that
 $A_i \star \lambda_y(\cS_{i}/\cS_{i-1}) - B_i$ may be pushed forward, and that
 \begin{gather}\label{eqn:nts}
 \pi_*\left(A_i \star \lambda_y(\cS_{i}/\cS_{i-1}) - B_i\right) = \lambda_y(\cS_{i+1}/\cS_{i-1}) .
 \end{gather}
 The hypothesis $r_i - r_{i-1} =1$ implies that $\cS_i/\cS_{i-1}$ is a line bundle, and that
 \begin{gather*}%\label{E:int1}
 A_i \star \lambda_y(\cS_{i}/\cS_{i-1}) - B_i\\
 \qquad{} = \lambda_y(\cS_{i+1}/\cS_{i}) \star \lambda_y(\cS_{i}/\cS_{i-1})
 + y^{r_{i+1}-r_{i-1}} \frac{Q_i}{1-Q_i} \det (\cS_{i+1}/\cS_{i}) \star \det (\cS_{i}/\cS_{i-1}) .
 \end{gather*}
 By \eqref{eqn:pushforwards-lambda}, we have
 \[
 \pi_*(\lambda_y(\cS_{i+1}/\cS_{i}))=\lambda_y(\cS_{i+1}/\cS_{i-1})-y^{r_i-r_{i-1}}\det(\cS_{i+1}/\cS_{i-1}), \qquad\pi_*(\lambda_y(\cS_{i}/\cS_{i-1}))=1.
 \]
 By \eqref{eqn:relX}, we have
 \[
 \frac{Q_i}{1-Q_i} \det (\cS_{i+1}/\cS_{i}) \star \det (\cS_{i}/\cS_{i-1}) = Q_i(1-Q_{i-1})\frac{\det\cS_{i+1}}{\det \cS_{i-1}}.
 \]
 As in the proof of (c), this can be pushed forward and its pushforward is $\det(\cS_{i+1}/\cS_{i-1})$. Putting these together, we have established \eqref{eqn:nts}.
\end{proof}

Recall that $Y^{(j)}=\bigl(Y^{(j)}_1,\dots,
{ Y^{(j)}_{r_{j+1}-r_j}}\bigr)$, $0\leq j\leq k$, are formal variables, $e_\ell$ denotes the $\ell$-th elementary symmetric polynomial, and $T_1,\dots, T_n \in \K_T(\pt)$ are given by
the decomposition of~$\C^n$ into one dimensional $T$-modules, that is, $\wedge^\ell(\C^n)=e_\ell(T_1,\dots,T_n)$.

\begin{Theorem}\label{thm:main}
 The ring $\QK_T(\Fl(\bfr,n))$ is isomorphic to $R\dbb{Q}/J_Q$, where
 \[
 R=\K_T(\pt)\bigl[e_1\bigl(Y^{(j)}\bigr),\dots,e_{r_{j+1}-r_j}\bigl(Y^{(j)}\bigr),\, 0\leq j\leq k\bigr],
 \]
 and $J_Q\subset R[\![Q]\!]=R[\![Q_1,\dots,Q_k]\!]$
is the ideal generated by the coefficients of $y$ in
\begin{align}\label{eqn:whit-short-gen}
 \prod_{\ell=1}^{n}(1+yT_\ell)-
 &\begin{vmatrix}
 A_0 & B_1 & & & \\
 1 & A_1 & B_2 & & \\ & \ddots & \ddots & \ddots & \\ & & 1 & A_{k-1} & B_{k}\\ & & & 1 & A_{k}
 \end{vmatrix},
 \end{align}
 where
 \[
 A_j=\prod_{\ell=1}^{r_{j+1}-r_j}\bigl(1+y Y^{(j)}_\ell\bigr)+B_j,\qquad B_j=y^{{r_{j+1}-r_j}}\frac{Q_j}{1-Q_j}\prod_{\ell=1}^{r_{j+1}-r_j}Y^{(j)}_\ell,
 \]
 with the convention that $Q_0=0$.

 More precisely,
% for $\mathbf{r}=(1,2,\dots,n-1)$,
there exists a $\K_T(\pt)[\![Q]\!]$-algebra isomorphism
\[
\Psi\colon\ R[\![Q]\!]/J_Q \to \QK_T(\Fl(r_1,\dots,r_k))
\]
that sends \smash{$e_\ell\bigl(Y^{(j)}\bigr)$} to $\wedge^\ell\br{\cS_{j+1}/\cS_j}$ for $j=0,\dots,k$ and $\ell=1,\dots,r_{j+1}-r_{j}$.
\end{Theorem}

\begin{proof}
It follows from Theorem~\ref{thm:whit-short} that $\Psi$ is a well-defined ring homomorphism.
To prove there are no other relations, we use~\cite[Theorem~4.1]{GMSXZZ:Nakayama}, which
states that a complete set of relations in the quantum (equivariant) K ring is
obtained by quantizing any complete set of relations in the ordinary (equivariant) K ring.
Therefore, we need to show that when one specializes each~$Q_i$ to~$0$, the resulting
ring is a presentation of $\K_T(\Fl(r_1,\dots,r_k))$.
% Geometrically,
The relations obtained this way
are the `Borel-type relations' of the $\lambda_y$ classes
\begin{equation}\label{E:pToda} \lambda_y(\cS_1) \cdot \lambda_y(\cS_2/\cS_1)\cdot \dots \cdot \lambda_y(\C^n/\cS_k) = \lambda_y(\C^n) . \end{equation}
Note that the relations~\eqref{E:pToda} can be obtained from the Whitney relations
\[ \lambda_y(\cS_i) \cdot \lambda(\cS_{i+1}/\cS_i) = \lambda_y(\cS_{i+1}) , \]
by eliminating the classes $\lambda_y(\cS_i)$ for $2 \le i \le k$.
(The quantization of this statement is done in the next section.)
Finally, it is known that
the Whitney relations form a full set of relations
in $\K_T(\Fl(r_1,\dots,r_k))$. This is essentially done by Lascoux~\cite[Section~7]{lascoux:anneau},
and we refer to~\cite[Proposition~5.1]{GMSXZZ:Nakayama} for a complete proof.
\end{proof}


% Original source lines 950--1026
\section{Whitney implies Toda}\label{sec:whitney}
In this section, we consider a different presentation of the quantum K ring, named the
{\em quantum~K Whitney presentation}.
This presentation quantizes
 relations $\lambda_y(\cS_i) \cdot \lambda_y(\cS_{i+1}/\cS_i) = \lambda_y(\cS_{i+1})$
 satisfied by the tautological subbundles in $\K_T(\Fl(\bfr,n))$. Informally, the
 Whitney presentation contains more (geometric) information than the Toda presentation,
 as it involves more generators, corresponding to the $\lambda_y$
 classes of the tautological subbundles, and their quotients. In~contrast, the Toda presentation
 only involves the quotient bundles.




The quantization was conjectured in~\cite{gu2023quantum,GMSXZZ:QKW},
generalizing the conjectures from
\cite{Gu:2020zpg} for Grassmannians. These conjectures have been proved in~\cite{gu2022quantum}
for Grassmannians, and in~\cite{GMSXZZ:QKW} for~$\Fl(1,n-1;n)$ case. The general case was recently announced in~\cite{huq2024quantum} using the abelian/non-abelian correspondence.
We note that the results in~\cite{huq2024quantum} are logically independent on those from~\cite{maeno2023presentation}, which were used to obtain the Toda presentation in the previous section.


Our main result of this section is that eliminating the additional variables of the Whitney presentation yields the Toda presentation. As an aside, we note that the proof of Theorem~\ref{thm:main} can be easily modified to show that the quantum K Whitney presentation of $\Fl(\bfr,n)$ follows from that of $\Fl(n)$.
We leave the details of this proof to the reader.

In what follows, $T$ can be a maximal torus in $\GL_n$.
Let
\[
 {X}^{(j)}=\bigl(X^{(j)}_1,\dots,X^{(j)}_{r_j}\bigr) \qquad\text{and}\qquad {Y}^{(j)}=\bigl(Y^{(j)}_1,\dots,
Y^{(j)}_{r_{j+1}-r_j}\bigr)
\]
denote formal variables for $j=1,\dots,k$ and denote by $X^{(k+1)}\coloneqq (T_1,\dots,T_n)$ the equivariant parameters in $\K_T(\pt)$. Let \smash{$e_\ell\bigl({X}^{(j)}\bigr)$} and \smash{$e_\ell\bigl({Y}^{(j)}\bigr)$} be the $\ell$-th elementary symmetric polynomials in ${X}^{(j)}$ and ${Y}^{(j)}$, respectively. Define the ring
 \[
 S=\K_T(\pt)\bigl[e_1\bigl(X^{(j)}\bigr),\dots, e_{r_j}\bigl(X^{(j)}\bigr),e_1\bigl(Y^{(j)}\bigr),\dots,e_{r_{j+1}-r_j}\bigl(Y^{(j)}\bigr)\bigr]_{j=1}^k ,
 \]
 and the ideal $I_Q\subset S[\![Q]\!]=S[\![Q_1,\dots,Q_k]\!]$
generated by the coefficients of $y$ in
\begin{gather}\label{eqn:qkrel}
 \prod_{\ell=1}^{r_j}\bigl(1+y X^{(j)}_\ell\bigr)\prod_{\ell=1}^{r_{j+1}-r_j}\bigl(1+y Y^{(j)}_\ell\bigr)-\prod_{\ell=1}^{r_{j+1}}\bigl(1+y X^{(j+1)}_\ell\bigr)\\
 \quad{}+y^{{r_{j+1}-r_j}}\frac{Q_j}{1-Q_j}\prod_{\ell=1}^{r_{j+1}-r_j}Y^{(j)}_\ell\Biggl(\prod_{\ell=1}^{r_j}\bigl(1+yX^{(j)}_\ell\bigr) -\prod_{\ell=1}^{r_{j-1}}\bigl(1+yX^{(j-1)}_\ell\bigr)\Biggr), \qquad j=1,\dots, k. \nonumber
\end{gather}
It was conjectured in~\cite{gu2023quantum,GMSXZZ:QKW} and proved in
\cite{huq2024quantum} that
there is an isomorphism of $\K_T(\pt)[\![Q]\!]$-algebras
 \begin{equation}\label{E:QKWhitney}
 \Phi\colon\ {S[\![Q]\!]}/I_Q\to \QK_T\br{\Fl(\bfr,n)}
 \end{equation}
 sending
 \[ e_\ell\bigl(X^{(j)}\bigr) \mapsto \wedge^\ell(\cS_{j}) \qquad \textrm{and} \qquad e_\ell\bigl(Y^{(j)}\bigr) \mapsto \wedge^\ell(\cS_{{j+1}}/\cS_{j}) . \]
 We refer to this as the (quantum K) {\em Whitney presentation}.

 \begin{Proposition}\label{prop:whit-toda}
 There is a natural isomorphism
 \[ S[\![Q]\!]/I_Q \simeq R\dbb{Q}/J_Q ,\]
 obtained by eliminating the indeterminates \smash{$X^{(j)}_\ell$}. In particular,
 the Whitney relations from~\eqref{eqn:qkrel} imply the
 Toda relations from~\eqref{eqn:whit-short-gen}.
 \end{Proposition}
 \begin{proof}
 Let
 \begin{align*}
 A_j=\prod_{\ell=1}^{r_{j+1}-r_j}\bigl(1+y Y^{(j)}_\ell\bigr)+B_j,\qquad B_j=y^{{r_{j+1}-r_j}}\frac{Q_j}{1-Q_j}\prod_{\ell=1}^{r_{j+1}-r_j}Y^{(j)}_\ell,
 \end{align*}
 so that~\eqref{eqn:qkrel} becomes
 \begin{align}\label{eqn:recur}
 A_j\prod_{\ell=1}^{r_j}\bigl(1+yX^{(j)}_\ell\bigr) -B_j\prod_{\ell=1}^{r_{j-1}}\bigl(1+yX^{(j-1)}_\ell\bigr) -\prod_{\ell=1}^{r_{j+1}}\bigl(1+y X^{(j+1)}_\ell\bigr).
 \end{align}
 Note that by Lemma~\ref{L:r-det}, relations given by~\eqref{eqn:recur} are equivalent to those given by
 \begin{align*}%\label{eqn:elim}
 \prod_{\ell=1}^{r_{j+1}}\bigl(1+y X^{(j+1)}_\ell\bigr)-
 &\begin{vmatrix}
 A_0 & B_1 & & & \\
 1 & A_1 & B_2 & & \\ & \ddots & \ddots & \ddots & \\ & & 1 & A_{j-1} & B_{j}\\ & & & 1 & A_{j}
 \end{vmatrix}\qquad \text{for $1\leq j\leq k$}.
 \end{align*}
 As a consequence, we can eliminate $e_1\bigl(X^{(j)}\bigr),\dots, e_{r_j}\bigl(X^{(j)}\bigr)$ for $2\leq j\leq k$, and be left with the relation~\eqref{eqn:whit-short-gen}.
 \end{proof}
\begin{Remark}\label{rmk:QKW} We note that our methods from the previous section can be adapted easily to show that $\Phi$ is an isomorphism for all partial flag varieties if and only if it is an isomorphism for $\Fl(n)$. \end{Remark}


% Original source lines 1385--1573
\appendix

\section[Toda relations from finite difference operators (after Anderson--Chen--Tseng)]{Toda relations from finite difference operators\\ (after Anderson--Chen--Tseng)}\label{sec:toda}

The proof of the Toda relations in~\cite{maeno.naito.sagaki:QKideal}
relies on Kato's earlier results~\cite{kato:loop}. For the quantum K~ring
$\QK_T(\Fl(n))$, there is another proof of these relations, using an argument
combining the results of Iritani, Milanov and
Tonita~\cite{Iritani:2013qka} with results of Givental
and Lee~\cite{givental_lee}. More precisely, it is shown in~\cite{Iritani:2013qka} that
the symbols of finite
difference operators annihilating the K-theoretic~$J$ function of a variety $X$ give relations in the quantum K ring of $X$. Givental and Lee's results from loc.\ cit.\ imply that the
K-theoretic~$J$ function of the complete flag variety is an eigenfunction of the
(finite difference) Toda Hamiltonians. This observation was made in the unpublished note~\cite{act:2017} of Anderson--Chen--Tseng, but removed from the published version of their paper. For the sake of completeness, we give a brief account below, and in the process fill in some of the details to make the argument complete.

We start with recalling the definition
of the K-theoretic $J$-function of the complete flag variety $X=\Fl(n)$.
Denote by $P_i= \wedge^i \cS_i$; it is known that these line bundles algebra generate~$\K_T(\Fl(n))$
over~$\K_T(\pt)$; see, e.g.,~\cite[Proposition~3.1]{GMSXZZ:Nakayama}. Furthermore,
the curve classes associated to the Novikov variables
$Q_i$ are dual to the classes
$c_1\bigl(\wedge^i \cS_i^*\bigr)$.
For a fixed effective (multi)degree~${d \in H_2(X)}$, let $L$ be the cotangent line bundle at the unique marked point on the moduli space $\overline{M}_{0,1}(X,d)$.
Let also $\phi^\alpha$, $\phi_\alpha$ denote Poincar\'e-dual bases for $\K_T(X)$. (For example,
one may take Schubert classes $\cO^w$, and their duals -- the ideal sheaves of the boundary of the opposite Schubert varieties.)
The small $J$-function of $X$, denoted by $J_X$, is defined by
\[
J_X(q)\coloneq (1-q)\prod_i P_i^{\frac{\ln(Q_i)}{\ln(q)}}\sum_{d,\alpha} Q^d\biggl\langle \frac{\phi_\alpha}{1-qL}\biggr\rangle_{0,1,d}\phi^\alpha .
\]
We will explain later the meaning and the effect of the factor
\smash{$P_i^{\ln(Q_i)/\ln(q)}$}. We also note that the presence
of the prefactors $(1-q)$
and \smash{$\prod_i P_i^{\ln(Q_i)/\ln(q)}$} varies in the literature.
Our description agrees with the one used by Givental and Lee in
\cite{givental_lee}, and corresponds to the function denoted by $\tilde{J}$ in~\cite{Iritani:2013qka}.

We recall some basics on the formalism of difference operators. Consider commuting
variables $q, x_1, \dots, x_n$, and define the difference operators
\[
T_i:= q^{x_i \partial_{x_i}} = \sum_{ k \ge 0}^\infty \frac{1}{k!}((\ln q)x_i \partial_{x_i} )^k.
\]
\big(More generally, for a differential operator $\mathfrak{f}$, one defines the $q$-difference operator $q^{\mathfrak{f}}={\rm e}^{(\ln q) \mathfrak{f}}=\sum_{j=0}^\infty \frac{1}{j!}((\ln q)\mathfrak{f})^j$.\big) Note that
\[
T_i \bigl(x_j^{\pm 1}\bigr) = \sum_{k=0}^\infty \frac{1}{k!} ( \ln(q) x_i \partial_{x_i} ) \bigl(x_j^{\pm 1}\bigr) = q^{\pm \delta_{ij}} x_j^{\pm 1} ,
\]
which explains the `difference operator' terminology. More generally, for any Laurent polynomial in commuting variables $x_i$, we have
\[
T_i f(x_1,\dots,x_i,\dots,x_n) = f(x_1,\dots,qx_i,\dots, x_n) ,
\]
i.e., $T_i$ is an automorphism of the Laurent polynomial ring
$\Z\bigl[q^{\pm 1}; x_1^{\pm 1}, \dots, x_{n}^{\pm 1}\bigr]$.
We use this expression to extend the definition of $T_i$ to any function in the indeterminates $q, x_1, \dots, x_n$.

Now consider the subring of Laurent polynomials
\[
\Z\bigl[q^{\pm 1}; Q_1^{\pm 1}, \dots, Q_{n-1}^{\pm 1}\bigr] \hookrightarrow \Z\bigl[q^{\pm 1}; x_1^{\pm 1}, \dots, x_{n}^{\pm 1}\bigr]
\]
obtained by sending \smash{$Q_i \mapsto q^{-1} \frac{x_{i+1}}{x_i}$}. The restriction of $T_i$ to this subring is given by
\begin{equation}\label{E: TiQ}
T_i = q^{-Q_{i}\partial_{Q_{i}}}q^{Q_{i-1}\partial_{Q_{i-1}}} ,
\end{equation}
where \smash{$q^{Q_i \partial_{Q_i}}$} are the difference operators on the subring in $Q_i$'s.

With that in mind, we can now explain the meaning of the factor \smash{$P^{\frac{\ln(Q_i)}{\ln(q)}}$}.
The difference operators \smash{$q^{Q_i \partial_{Q_i}}$} act on functions in $Q_i$'s, and
one calculates that
\[ q^{Q_i \partial_{Q_i}}\bigl(P^{\frac{\ln(Q_j)}{\ln(q)}}\bigr) =
P^{\frac{\ln(q^{\delta_{ij}} Q_j)}{\ln(q)}} = P^{\delta{ij}}
P^{\frac{\ln(Q_j)}{\ln(q)}} . \]
In other words, the factor \smash{$P^{\frac{\ln(Q_i)}{\ln(q)}}$} should be regarded as
a formal variable
which transforms according to the rule above under the difference operators.

The relations in the quantum K ring are given by the Hamiltonians of the finite difference
(or relativistic) Toda lattice. There is some ambiguity in the exact expressions for the
Toda Hamiltonians, since their construction depends on choices;
see, e.g.,~\cite[Remark~5]{givental_lee}. We follow here the approach from~\cite{gloI},
but we will also need to make some changes of variables,
in order to fit with the conventions in our main reference~\cite{givental_lee}.
For the convenience of the reader, we briefly included some of the details below.

The Hamiltonians of the $q$-deformed type $A$ Toda chain have the form
\begin{align}\label{eqn:hamil}
 H_k=\sum_{0= i_0<\dots<i_k\leq n} \prod_{l=1}^{k}\biggl(1-\frac{x_{i_l}}{x_{i_l-1}}\biggr)^{1-\delta_{i_l-i_{l-1},1}} \prod_{l=1}^{k} T_{i_l},\qquad k=1,\dots,n,
\end{align}
where $q$ and $x_i$ are commuting variables, and $T_i=q^{x_i\partial_{x_i}}$ is the $q$-difference operator above. It was proved in~\cite{gloI} that the operators $H_k$ are limits of Macdonald operators, and the latter are known to commute.
This implies that $H_k$ also commute.

As above, let $Q_i=q^{-1}{x_{i+1}}{x_i}^{-1}$, with $Q_0=Q_n=0$. Then, using
\eqref{E: TiQ}, one can rewrite~\eqref{eqn:hamil} as
\begin{align*}
 H_k=\sum_{0= i_0<\dots<i_k\leq n} \prod_{l=1}^{k}(1-qQ_{i_l-1})^{1-\delta_{i_l-i_{l-1},1}} \prod_{l=1}^{k} q^{-Q_{i_l}\partial_{Q_{i_l}}}q^{Q_{i_l-1}\partial_{Q_{i_l-1}}},\qquad k=1,\dots,n.
\end{align*}
Replacing $q$ by $q^{-1}$, we obtain
\begin{align*}%\label{E:Toda-ops-changed}
 \widehat{H}_k&{}= \sum_{0= i_0<\dots<i_k\leq n}\prod_{l=1}^{k}\bigl(1-q^{-1}Q_{{i_l}-1}\bigr)^{1-\delta_{i_l-i_{l-1},1}} \prod_{l=1}^{k} q^{Q_{i_l}\partial_{Q_{i_l}}-Q_{i_l-1}\partial_{Q_{{i_l}-1}}}\\
 &{}=\sum_{0= i_0<\dots<i_k\leq n} \prod_{l=1}^{k} q^{Q_{i_l}\partial_{Q_{i_l}}-Q_{i_l-1}\partial_{Q_{{i_l}-1}}}\prod_{l=1}^{k}(1-Q_{{i_l}-1})^{1-\delta_{i_l-i_{l-1},1}}
 ,\qquad k=1,\dots,n .
 \end{align*}
\begin{Remark} The substitutions above
ensure that the first Hamiltonian $\widehat{H}_1$ agrees
with the one used in~\cite{givental_lee}. The substitution chosen in~\cite{act:2017}
produces similar
operators, but with the $q$-shifts and the Novikov terms in the opposite order.
\end{Remark}
The following key result of Givental and Lee~\cite[Theorem~2]{givental_lee} shows that the $J$ function is an eigenfunction for $J_{\Fl(n)}$.

\begin{Theorem}[Givental--Lee]\label{thm:GLeemain}
 $\widehat{H}_1J_{\Fl(n)}=\mathbb{C}^n J_{\Fl(n)}$.
\end{Theorem}

We also need the following lemma of Givental--Lee~\cite{givental_lee}.

\begin{Lemma}[{\cite[p.~9]{givental_lee}}]\label{lemma:modQ}
Let $D$ be a difference operator commuting with \smash{$\widehat{H}_1$}. Then, if $J$ is an eigenfunction of $D$ modulo $Q$, then $J$ is an eigenfunction of $D$ whose eigenvalue is the same as the one modulo $Q$.
\end{Lemma}

From this, we deduce that $J_{\Fl(n)}$ is an eigenfunction of the higher Toda Hamiltonians, using their commutativity with \smash{$\widehat{H}_1$} and by computing their eigenvalues modulo $Q$.

\begin{Corollary}\label{cor:Hk-eigenval} For any $1 \le k \le n$, the following holds:
\[ \widehat{H}_kJ_{\Fl(n)}=\wedge^k(\mathbb{C}^n) J_{\Fl(n)} . \]
\end{Corollary}
\begin{proof} The case $k=1$ is Theorem~\ref{thm:GLeemain}. Suppose $2 \le k \le n$. Since \smash{$\widehat{H}_k$} commutes with \smash{$\widehat{H}_1$}, we need only verify that $J_{\Fl(n)}$ is an eigenfunction of \smash{$\widehat{H}_k$} modulo $Q$, thanks to Lemma~\ref{lemma:modQ}.

To this end, we first observe
\begin{align*}
\widehat{H}_k J_{\Fl(n)}&{}=\widehat{H}_k\biggl((1-q)\prod_i P_i^{\frac{ln(Q_i)}{ln(q)}}\biggr) + o(Q_i) \\
&{}=
\sum_{0= i_0<\dots<i_k\leq n}\prod_{l=1}^k \frac{P_{i_l}}{P_{i_l-1}}\biggl((1-q)\prod_i P_i^{\frac{ln(Q_i)}{ln(q)}}\biggr)+o(Q_i) .
\end{align*}
Thus, modulo $Q$, we have the eigenvalue equation
\begin{align*}
\widehat{H}_k J_{\Fl(n)}&{}=\sum_{0= i_0<\dots<i_k\leq n}
\prod_{l=1}^k \frac{P_{i_l}}{P_{i_l-1}}J_{\Fl(n)}\\
&{}=
e_k\biggl(\frac{P_1}{P_0},\frac{P_2}{P_1},\dots,\frac{P_{n}}{P_{n-1}}\biggr)J_{\Fl(n)}=\wedge^k(\mathbb{C}^n)J_{\Fl(n)} .
\tag*{\qed}
\end{align*}
\renewcommand{\qed}{}
\end{proof}

We now use~\cite[Proposition~2.12]{Iritani:2013qka} which shows that the symbols of
Toda Hamiltonians give relations in quantum K theory.

\begin{Theorem}[Iritani--Milanov--Tonita]\label{thm:IMT}
Let $D=D\bigl(q^{Q_i\partial_{Q_i}},q,Q,\Lambda_i\bigr)$ be any $q$-difference operator
with coefficients in $\K_T(\pt)[q^{\pm 1}][\![Q_i]\!]$, such that it is regular
at $q=1$. Then
\[DJ_X=0\implies D\bigl(\widehat{P}_i,1,Q,\Lambda_i\bigr)=0 \in \QK_T(X) .\]
\end{Theorem}

\begin{Remark}
The result of Iritani, Milanov and Tonita is stated non-equivariantly, and for the
{\em big} quantum K ring and the corresponding big $J$ function. However, an inspection of their proof shows that it works in the equivariant situation as well.
Furthermore, if one starts with the small quantum K ring, then all arguments extend
to that situation, and the result also holds for the small quantum K ring and the small $J$ function. For further details, see~\cite{huq2024relations}.
\end{Remark}
One subtle point is that $\widehat{P}_i$ is a certain $Q$-deformation
of the line bundle $P_i$: it is the restriction to the small
quantum K ring
of an operator denoted by
$A_{i,{\rm com}}$ in~\cite[Corollary~2.9]{Iritani:2013qka}, which
arises as a solution to a certain Lax-type equation.
However, results of both Anderson, Chen, Tseng, and Iritani
in~\cite[Lemma~6]{anderson2022finite}, and also
by Kato in~\cite[Theorem~1.35]{kato:loop}
show that in fact no quantization is needed.

\begin{Proposition}\label{thm:ACTI-Kato} For the flag variety $\Fl(n)$,
$\widehat{\det(\mathcal{S}_i)}=\det(\mathcal{S}_i)$.
\end{Proposition}
We note in passing that an analogue of Proposition~\ref{thm:ACTI-Kato} holds for any homogeneous
space~$G/P$, but we do not need this generality here.

Combining Theorem~\ref{thm:IMT} and Proposition~\ref{thm:ACTI-Kato} with Corollary~\ref{cor:Hk-eigenval}
yields the following corollary.
\begin{Corollary}%\label{cor:QKToda}
The following identities hold
in $\QK_T(\Fl(n))$:
\begin{equation*}
 \sum_{0= i_0<\dots<i_k\leq n}\ \prod_{l=1}^{k}\frac{P_{i_l}}{P_{i_{l}-1}}(1-Q_{i_l-1})^{1-\delta_{i_l-i_{l-1},1}}=\wedge^k \C^n,\qquad k=1,\dots,n,
\end{equation*}
where $P_0=P_n=1$.
\end{Corollary}

\begin{Remark}
Theorem 4.9 of~\cite{koroteev} gives a presentation of the quasimap quantum $K$-ring of $T^*Fl$ whose limit to the is described in Theorem 5.5. The relations are based on the trigonometric Ruijsenaars--Schneider model. After further taking into account a restriction from $\GL_n$ to $\SL_n$, the `Toda limit' recovers the relations in this paper. We are grateful to P. Koroteev who explained this procedure to us.
\end{Remark}

\begin{thebibliography}{99}
\bibitem[1]{act:2017}
Anderson D., Chen L., Tseng H.-H., On the quantum {K}-ring of the flag manifold,
 \href{http://arxiv.org/abs/1711.08414}{arXiv:1711.08414}.

\bibitem[2]{anderson2022finite}
Anderson D., Chen L., Tseng H.-H., On the finiteness of quantum {K}-theory of
 a~homogeneous space (with an appendix by {H}iroshi {I}ritani),
 \href{https://doi.org/10.1093/imrn/rnaa108}{\textit{Int. Math. Res. Not.}}
 \textbf{2022} (2022), 1313--1349,
 \href{http://arxiv.org/abs/1804.04579}{arXiv:1804.04579}.

\bibitem[5]{BCMP:qkpos}
Buch A.S., Chaput P.-E., Mihalcea L.C., Perrin N., Positivity in minuscule
 quantum~{K} theory, \href{http://arxiv.org/abs/2205.08630}{arXiv:2205.08630}.

\bibitem[6]{BCMP:qkchev}
Buch A.S., Chaput P.-E., Mihalcea L.C., Perrin N., A~{C}hevalley formula for the
 equivariant quantum {$K$}-theory of cominuscule varieties,
 \href{https://doi.org/10.14231/ag-2018-015}{\textit{Algebr. Geom.}}
 \textbf{5} (2018), 568--595,
 \href{http://arxiv.org/abs/1604.07500}{arXiv:1604.07500}.

\bibitem[7]{buch.m:qk}
Buch A.S., Mihalcea L.C., Quantum {$K$}-theory of {G}rassmannians,
 \href{https://doi.org/10.1215/00127094-2010-218}{\textit{Duke Math.~J.}}
 \textbf{156} (2011), 501--538.

\bibitem[8]{chaput.perrin:rationality}
Chaput P.-E., Perrin N., Rationality of some {G}romov--{W}itten varieties and
 application to quantum {$K$}-theory,
 \href{https://doi.org/10.1142/S0219199711004166}{\textit{Commun. Contemp.
 Math.}} \textbf{13} (2011), 67--90,
 \href{http://arxiv.org/abs/0905.4394}{arXiv:0905.4394}.

\bibitem[10]{etingof}
Etingof P., Whittaker functions on quantum groups and {$q$}-deformed {T}oda
 operators, in Differential Topology, Infinite-dimensional {L}ie Algebras, and
 Applications, \textit{Amer. Math. Soc. Transl. Ser.~2}, Vol. 194,
 \href{https://doi.org/10.1090/trans2/194/02}{American Mathematical Society},
 Providence, RI, 1999, 9--25,
 \href{http://arxiv.org/abs/math.QA/9901053}{arXiv:math.QA/9901053}.

\bibitem[13]{gloI}
Gerasimov A., Lebedev D., Oblezin S., On {$q$}-deformed
 {${\mathfrak{gl}}_{l+1}$}-{W}hittaker function.~{I},
 \href{https://doi.org/10.1007/s00220-009-0917-y}{\textit{Comm. Math. Phys.}}
 \textbf{294} (2010), 97--119,
 \href{http://arxiv.org/abs/0805.3754}{arXiv:0805.3754}.

\bibitem[14]{givental:onwdvv}
Givental A., On the {WDVV} equation in quantum {$K$}-theory,
 \href{https://doi.org/10.1307/mmj/1030132720}{\textit{Michigan Math.~J.}}
 \textbf{48} (2000), 295--304,
 \href{http://arxiv.org/abs/math.AG/0003158}{arXiv:math.AG/0003158}.

\bibitem[15]{givental_lee}
Givental A., Lee Y.-P., Quantum {$K$}-theory on flag manifolds,
 finite-difference {T}oda lattices and quantum groups,
 \href{https://doi.org/10.1007/s00222-002-0250-y}{\textit{Invent. Math.}}
 \textbf{151} (2003), 193--219,
 \href{http://arxiv.org/abs/math.AG/0108105}{arXiv:math.AG/0108105}.

\bibitem[16]{Gorbounov:2014}
Gorbounov V., Korff C., Quantum integrability and generalised quantum
 {S}chubert calculus,
 \href{https://doi.org/10.1016/j.aim.2017.03.030}{\textit{Adv. Math.}}
 \textbf{313} (2017), 282--356,
 \href{http://arxiv.org/abs/1408.4718}{arXiv:1408.4718}.

\bibitem[17]{gu2023quantum}
Gu W., Mihalcea L., Sharpe E., Xu W., Zhang H., Zou H., Quantum~{K} theory
 rings of partial flag manifolds,
 \href{https://doi.org/10.1016/j.geomphys.2024.105127}{\textit{J.~Geom.
 Phys.}} \textbf{198} (2024), 105127, 30~pages,
 \href{http://arxiv.org/abs/2306.11094}{arXiv:2306.11094}.

\bibitem[18]{Gu:2020zpg}
Gu W., Mihalcea L., Sharpe E., Zou H., Quantum~{K} theory of symplectic
 {G}rassmannians,
 \href{https://doi.org/10.1016/j.geomphys.2022.104548}{\textit{J.~Geom.
 Phys.}} \textbf{177} (2022), 104548, 38~pages,
 \href{http://arxiv.org/abs/2008.04909}{arXiv:2008.04909}.

\bibitem[19]{GMSXZZ:Nakayama}
Gu W., Mihalcea L.C., Sharpe E., Xu W., Zhang H., Zou H., A~{N}akayama result
 for the quantum~{K} theory of homogeneous spaces, \textit{\'Epijournal
 G\'eom. Alg\'ebrique}, {t}o appear,
 \href{http://arxiv.org/abs/2507.15183}{arXiv:2507.15183}.

\bibitem[20]{GMSXZZ:QKW}
Gu W., Mihalcea L.C., Sharpe E., Xu W., Zhang H., Zou H., Quantum~{K} {W}hitney
 relations for partial flag varieties,
 \href{http://arxiv.org/abs/2310.03826}{arXiv:2310.03826}.

\bibitem[21]{gu2022quantum}
Gu W., Mihalcea L.C., Sharpe E., Zou H., Quantum~{K} theory of {G}rassmannians,
 {W}ilson line operators and {S}chur bundles,
 \href{https://doi.org/10.1017/fms.2025.10088}{\textit{Forum Math. Sigma}}
 \textbf{13} (2025), e140, 38~pages,
 \href{http://arxiv.org/abs/2208.01091}{arXiv:2208.01091}.

\bibitem[22]{huq2024quantum}
Huq-Kuruvilla I., Quantum {K}-Rings of partial flag varieties, {C}oulomb
 branches, and the {B}ethe ansatz,
 \href{http://arxiv.org/abs/2409.15575}{arXiv:2409.15575}.

\bibitem[23]{huq2024relations}
Huq-Kuruvilla I., Relations in twisted quantum {K}-rings,
 \href{http://arxiv.org/abs/2406.00916}{arXiv:2406.00916}.

\bibitem[26]{Iritani:2013qka}
Iritani H., Milanov T., Tonita V., Reconstruction and convergence in quantum
 {$K$}-theory via difference equations,
 \href{https://doi.org/10.1093/imrn/rnu026}{\textit{Int. Math. Res. Not.}}
 \textbf{2015} (2015), 2887--2937,
 \href{http://arxiv.org/abs/1309.3750}{arXiv:1309.3750}.

\bibitem[27]{kapranov:Gr}
Kapranov M.M., On~the derived category of coherent sheaves on Grassmann
 manifolds,
 \href{https://doi.org/10.1070/IM1985v024n01ABEH00122}{\textit{Math.
 USSR-Izv.}} \textbf{24} (1985), 183--192.

\bibitem[28]{kato2019quantum}
Kato S., On quantum $K$-groups of partial flag manifolds,
 \href{http://arxiv.org/abs/1906.09343}{arXiv:1906.09343}.

\bibitem[29]{kato:loop}
Kato S., Loop structure on equivariant {$K$}-theory of semi-infinite flag
 manifolds, \href{https://doi.org/10.4007/annals.2025.202.3.2}{\textit{Ann. of
 Math.}} \textbf{202} (2025), 1001--1075,
 \href{http://arxiv.org/abs/1805.01718}{arXiv:1805.01718}.

\bibitem[30]{koroteev}
Koroteev P., Pushkar P.P., Smirnov A.V., Zeitlin A.M., Quantum {K}-theory of
 quiver varieties and many-body systems,
 \href{https://doi.org/10.1007/s00029-021-00698-3}{\textit{Selecta
 Math.~(N.S.)}} \textbf{27} (2021), 87, 40~pages,
 \href{http://arxiv.org/abs/1705.10419}{arXiv:1705.10419}.

\bibitem[34]{lascoux:anneau}
Lascoux A., Anneau de {G}rothendieck de la vari\'et\'e de drapeaux, in The
 {G}rothendieck {F}estschrift, {V}ol.~{III}, \textit{Progr. Math.}, Vol.~88,
 \href{https://doi.org/10.1007/978-0-8176-4576-2_1}{Birkh\"auser}, Boston, MA,
 1990, 1--34.

\bibitem[35]{lee:QK}
Lee Y.-P., Quantum {$K$}-theory.~{I}. {F}oundations,
 \href{https://doi.org/10.1215/S0012-7094-04-12131-1}{\textit{Duke Math.~J.}}
 \textbf{121} (2004), 389--424,
 \href{http://arxiv.org/abs/math/0105014}{arXiv:math/0105014}.

\bibitem[38]{LNSX:QKdiv}
Lenart C., Naito S., Sagaki D., Mihalcea L.C., Xu W., Quantum {K}-theoretic
 divisor axiom for flag manifolds,
 \href{http://arxiv.org/abs/2505.16150}{arXiv:2505.16150}.

\bibitem[39]{maeno.naito.sagaki:QKideal}
Maeno T., Naito S., Sagaki D., A~presentation of the torus-equivariant quantum
 {$K$}-theory ring of flag manifolds of type~{$A$}, {P}art~{I}: {T}he defining
 ideal, \href{https://doi.org/10.1112/jlms.70095}{\textit{J.~Lond. Math.
 Soc.}} \textbf{111} (2025), e70095, 43~pages,
 \href{http://arxiv.org/abs/2302.09485}{arXiv:2302.09485}.

\bibitem[40]{maeno2023presentation}
Maeno T., Naito S., Sagaki D., A~presentation of the torus-equivariant quantum
 {$K$}-theory ring of flag manifolds of type~{$A$}, {P}art~{II}: quantum
 double {G}rothendieck polynomials,
 \href{https://doi.org/10.1017/fms.2024.147}{\textit{Forum Math. Sigma}}
 \textbf{13} (2025), e19, 26~pages,
 \href{http://arxiv.org/abs/2305.17685}{arXiv:2305.17685}.

\bibitem[43]{sinha2024quantum}
Sinha S., Zhang M., Quantum $K$-invariants via {Q}uot schemes~{I},
 \href{http://arxiv.org/abs/2406.12191}{arXiv:2406.12191}.

\end{thebibliography}
\end{document}
